\subsection{点处的可微函数}

1.2.1. 设 \(f\) 在 E 中点 \(x_{0}\) 的某邻域上定义且取值于 F。若存在从 E 到 F 的连续仿射函数 \(f\)，使得 \(x_{0}\) 与 \(v\) 在 \(x_{0}\) 处的接触阶 \(\geqslant 1\)，则称 \(f\) 在 x＿0 处可微。此映射 \(v\) 唯一；存在唯一的从 E 到 F 的连续线性映射，记作 \(Df(x_{0})\)，使得：

\[
v (x) = v \left(x _ {0}\right) + D f \left(x _ {0}\right). \left(x - x _ {0}\right).
\]

若在 E 上选取一个范数，则这等价于：

\(f(x_0 + h) \equiv f(x_0) + Df(x_0).h \mod o(\|h\|)\)（其中 \(h\) 趋于 0），

这也可写成如下形式：

\[
\lim _ {h \rightarrow 0, h \neq 0} \frac {\| f (x _ {0} + h) - f (x _ {0}) - D f (x _ {0}) . h \| _ {\gamma}}{\| h \|} = 0
\]

对 F 上每个连续半范数 γ。

称 \(Df(x_{0})\) 中的元素 \(\mathcal{L}(\mathrm{E},\mathrm{F})\) 为 f 在 \(x_{0}\) 处的导数。有时用 \(D_{h}f(x_{0})\) 表示 \(Df(x_{0}).h\)。它是 F 中由下式定义的元素：

\[
\mathrm{D} _ {h} f (x _ {0}) = \lim _ {t \rightarrow 0, t \neq 0} \frac {f (x _ {0} + t h) - f (x _ {0})}{t}.
\]

1.2.2. 若函数 \(f\) 在 \(x_{0}\) 处可微，且对每个定义 E 的拓扑的范数都有下列关系，则称 f 在 \(x_{0}\) 处严格可微：

\[
f (y) - f (z) \equiv D f (x _ {0}). (y - z) \quad \text {   mod   } o (\| y - z \|)
\]

当 \((y, z)\) 在 \((x_0, x_0)\) 中趋于 \(E \times E\) 时成立。为此，只需对定义 \(E\) 的拓扑的某个范数满足此条件即可。此外，若 \(E\) 和 \(F\) 都是赋范空间，则对每个数 \(c > \|Df(x_0)\|\)，存在 \(x_{0}\) 的一个邻域 V，使得对 V 中的 y、z 有 \(\|f(y)-f(z)\|\leqslant c.\|y-z\|\)；这蕴含 f 在 V 中一致连续。

1.2.3. 函数 \(f\) 在 \(x_{0}\) 处可微或严格可微这一性质只取决于 \(f\) 在 \(x_{0}\) 处的芽。在 \(x_{0}\) 处可微的函数芽构成所有芽空间的向量子空间 \(\mathcal{V}\)，且从 \(f\mapsto\mathrm{D}f(x_{0})\) 到 \(\mathcal{V}\) 的映射 \(\mathcal{L}(\mathrm{E};\mathrm{F})\) 是线性的。在 \(x_{0}\) 处严格可微的函数芽构成 \(\mathcal{V}\) 的向量子空间。

1.2.4. 在 \(x_0\) 处可微的函数在 \(x_0\) 处连续。

1.2.5. 当 E = K 时，映射 \(u \mapsto u(1)\) 是从 \(\mathcal{L}(\mathbf{E}; \mathbf{F})\) 到 F 的同构；若函数 f 在 \(x_{0}\) 处可微，则元素

\[
f ^ {\prime} (x _ {0}) = \mathrm{D} f (x _ {0}). 1
\]

正是 Fonct. Var. Réelle，第 I 章，第 1 节，第 6 号注 2 中所给意义下的 \(f\) 在 \(x_0\) 处的导数。

